Это один из \textbf{двух альтернативных вариантов} лабораторной работы
№4. Выберите и защитите \textbf{один}: базовый вариант 1 --- программа
факториала на C с GCC и GDB; вариант 2 --- загрузочный сектор и журнал
оценок на NASM в QEMU. Выполнять оба варианта не требуется. Ниже
приведена полная методичка варианта 2.

Рекомендуется \textbf{Debian Linux актуальной версии} --- установленный
на компьютере или в виртуальной машине. Можно использовать и другой
дистрибутив Linux, если доступны перечисленные ниже инструменты. В любом
случае нужны графический рабочий стол и право устанавливать пакеты.
Домашний опыт с \texttt{memory.c} рассчитан на Linux x86-64 (модель
LP64, порядок байтов little-endian): на компьютере с другой архитектурой
используйте совместимую виртуальную машину x86-64. Если Linux пока не
установлен, используйте время семинара №4 для установки системы или
настройки виртуальной машины. Задания, которые не успели выполнить в
аудитории, доделайте дома.

Распакуйте архив материалов семинара: пять исходных файлов и короткий
\texttt{README.md} должны лежать рядом. Команды подготовки ниже
выполняйте из этой папки, остальные --- из созданного в ней каталога
\texttt{work/}. Выполняйте команды самостоятельно, проверяйте экран,
вывод терминала и байты образов; заносите ожидаемое и наблюдаемое в
отчёт.

\hypertarget{ux43fux43eux440ux44fux434ux43eux43a-ux441ux434ux430ux447ux438-ux43bux430ux431ux43eux440ux430ux442ux43eux440ux43dux44bux445-ux440ux430ux431ux43eux442}{%
\subsubsection{Порядок сдачи лабораторных
работ}\label{ux43fux43eux440ux44fux434ux43eux43a-ux441ux434ux430ux447ux438-ux43bux430ux431ux43eux440ux430ux442ux43eux440ux43dux44bux445-ux440ux430ux431ux43eux442}}

\begin{enumerate}
\def\labelenumi{\arabic{enumi}.}
\tightlist
\item
  \textbf{На занятии} выполните часть работы, запишите команды, расчёты
  и наблюдаемый результат и покажите его преподавателю. Если в процессе
  работы возникли вопросы, задайте их преподавателю.
\item
  \textbf{К следующей лабораторной} завершите работу и подготовьте
  отчёт: что сделали, как проверили, что ожидали и что наблюдали.
  Изучите использованные команды и понятия; подготовьтесь объяснить,
  как работает ваша программа, применённые механизмы и устройства компьютера.
  Сохраните исходники и результаты, чтобы повторить демонстрацию.
  \textbf{Отправьте преподавателю отчёт в формате PDF или HTML.}
\item
  \textbf{Во время следующей лабораторной} индивидуально защитите
  предыдущую работу: покажите преподавателю отчёт и работающий
  результат, объясните свои действия и ответьте на вопросы. Преподаватель
  сообщит, если нужно что-то исправить.
\item
  \textbf{На последнем занятии} защитите предыдущую лабораторную работу
  и сдайте оставшиеся задолженности. Преподаватель объявит текущую
  накопленную оценку за лабораторные работы.
\end{enumerate}

В конце семинара №4 покажите преподавателю, что успели сделать:
результат установки Debian или настройки виртуальной машины,
\texttt{HELLO\ BIOS} в QEMU либо работающий журнал. Запишите выполненные
шаги и то, что наблюдали. К семинару №5 завершите задания и отчёт по
форме в конце методички; сохраните изменённый \texttt{journal.asm},
необходимые для повторного запуска образы и результаты опыта с C/GDB.
Во время семинара №5 покажите отчёт
и результат преподавателю, повторите ключевую проверку и объясните свои
изменения.

\hypertarget{ux443ux441ux442ux430ux43dux43eux432ux438ux442ux435-ux438ux43dux441ux442ux440ux443ux43cux435ux43dux442ux44b}{%
\subsubsection{Установите
инструменты}\label{ux443ux441ux442ux430ux43dux43eux432ux438ux442ux435-ux438ux43dux441ux442ux440ux443ux43cux435ux43dux442ux44b}}

Для работы понадобятся:

\begin{itemize}
\tightlist
\item
  QEMU с эмулятором ПК \texttt{qemu-system-i386} и графическим окном для
  гостевой машины;
\item
  ассемблер NASM и входящая в его пакет утилита \texttt{ndisasm};
\item
  компилятор GCC для C и отладчик GDB;
\item
  binutils (в том числе \texttt{objdump}) и утилита \texttt{file} для
  определения типа файла;
\item
  обычные команды оболочки и файловые утилиты: \texttt{cp},
  \texttt{mkdir}, \texttt{wc}, \texttt{od}, \texttt{dd},
  \texttt{truncate}, \texttt{diff}, \texttt{printf}.
\end{itemize}

Следующие команды установки приведены \textbf{для Debian Linux}:

\begin{verbatim}
sudo apt update
sudo apt install qemu-system-x86 qemu-system-gui \
  nasm gcc libc6-dev gdb binutils file coreutils diffutils
nasm -v
qemu-system-i386 --version
gcc --version
gdb --version
objdump --version
\end{verbatim}

\texttt{apt\ update} обновляет список пакетов; \texttt{apt\ install}
устанавливает перечисленные инструменты, заголовки и библиотеки для
компиляции C (\texttt{libc6-dev}), а также базовые файловые утилиты.
Пакет \texttt{qemu-system-gui} добавляет окно GTK; \texttt{file}
определяет тип файла по содержимому. Команды с
\texttt{-\/-version}/\texttt{-v} позволяют проверить наличие
инструментов; версии пакетов могут отличаться.

Если вы выбрали другой дистрибутив Linux, названия пакетов, команды
установки и настройка графического запуска QEMU могут быть другими.
Подбирайте их и разбирайтесь с возникающими особенностями
самостоятельно: преподаватели ориентируются на Debian Linux и не могут
детально изучить каждый дистрибутив, который студенты установили
индивидуально.

\hypertarget{ux441ux43eux431ux435ux440ux438ux442ux435-ux441ux432ux43eux44e-ux43fux435ux440ux432ux443ux44e-ux43fux440ux43eux433ux440ux430ux43cux43cux443}{%
\subsubsection{Соберите свою первую
программу}\label{ux441ux43eux431ux435ux440ux438ux442ux435-ux441ux432ux43eux44e-ux43fux435ux440ux432ux443ux44e-ux43fux440ux43eux433ux440ux430ux43cux43cux443}}

Из \textbf{папки распакованных материалов} сделайте собственную копию
исходника и перейдите в рабочий каталог. \texttt{cp\ -n} не заменяет уже
существующий файл --- ваши правки останутся на месте при повторном
выполнении:

\begin{verbatim}
mkdir -p work
cp -n hello.asm work/hello.asm
cd work
\end{verbatim}

Прочитайте \texttt{hello.asm}: \texttt{bits\ 16}, \texttt{org\ 0x7c00},
установка \texttt{DS} и \texttt{SS:SP}, вывод через \texttt{INT\ 10h},
нулевой байт после строки, \texttt{times} и \texttt{dw\ 0xaa55}.
Соберите и исследуйте загрузочный сектор:

\begin{verbatim}
nasm -f bin hello.asm -o hello.bin
wc -c hello.bin
od -An -tx1 -j 510 -N 2 hello.bin
ndisasm -b 16 hello.bin
cp hello.bin hello.img
truncate -s 1440K hello.img
wc -c hello.img
\end{verbatim}

\begin{longtable}[]{@{}
  >{\raggedright\arraybackslash}p{(\columnwidth - 2\tabcolsep) * \real{0.2700}}
  >{\raggedright\arraybackslash}p{(\columnwidth - 2\tabcolsep) * \real{0.6900}}@{}}
\toprule
\begin{minipage}[b]{\linewidth}\raggedright
Параметр
\end{minipage} & \begin{minipage}[b]{\linewidth}\raggedright
Что означает
\end{minipage} \\
\midrule
\endhead
\texttt{-f\ bin} & Плоские машинные байты без заголовка ELF \\
\texttt{-o\ hello.bin} & Имя результата NASM \\
\texttt{wc\ -c} & Число байтов файла: ожидайте \textbf{512} для
сектора \\
\texttt{od\ -An\ -tx1} & Вывод байтов hex, без колонки адресов \\
\texttt{-j\ 510\ -N\ 2} & Пропустить 510 байтов и прочитать два: подпись
\textbf{55 AA} \\
\texttt{ndisasm\ -b\ 16} & Показать команды, интерпретируя их как
16-битный код \\
\texttt{truncate\ -s\ 1440K} & Установить размер дискеты: 1440×1024
байта \\
\bottomrule
\end{longtable}

\texttt{ndisasm} не знает, где закончился выполняемый код: строку,
дополнение нулями и подпись он тоже попытается показать как инструкции.
Сопоставляйте с исходником прежде всего первые команды.

Запустите эмулируемый компьютер:

\begin{verbatim}
qemu-system-i386 \
  -machine pc -m 16M -boot order=a \
  -drive file=hello.img,format=raw,if=floppy \
  -display gtk -no-reboot
\end{verbatim}

\begin{longtable}[]{@{}
  >{\raggedright\arraybackslash}p{(\columnwidth - 2\tabcolsep) * \real{0.2700}}
  >{\raggedright\arraybackslash}p{(\columnwidth - 2\tabcolsep) * \real{0.6900}}@{}}
\toprule
\begin{minipage}[b]{\linewidth}\raggedright
Флаг QEMU
\end{minipage} & \begin{minipage}[b]{\linewidth}\raggedright
Значение
\end{minipage} \\
\midrule
\endhead
\texttt{-machine\ pc} & BIOS-совместимый виртуальный ПК \\
\texttt{-m\ 16M} & 16 MiB оперативной памяти гостя \\
\texttt{-boot\ order=a} & Сначала загрузочная дискета A: \\
\texttt{if=floppy} & Образ подключён как дискета; \texttt{file} задаёт
путь, \texttt{format=raw} --- плоский формат \\
\texttt{-display\ gtk} & Показать отдельное окно гостевой машины \\
\texttt{-no-reboot} & Не перезапускать гостя автоматически после
сброса \\
\bottomrule
\end{longtable}

\textbf{Наблюдение:} появилось \texttt{HELLO\ BIOS}. Поменяйте строку в
\texttt{hello.asm} на своё короткое латинское сообщение, повторите
команды начиная с \texttt{nasm} и убедитесь, что на экране уже новый
текст. Остановите QEMU закрытием окна или \texttt{Ctrl+C}.
\texttt{hello.bin} --- код для BIOS, не программа Linux для запуска
через \texttt{./hello.bin}.

\hypertarget{ux441ux43eux431ux435ux440ux438ux442ux435-ux433ux43eux442ux43eux432ux44bux439-ux436ux443ux440ux43dux430ux43b}{%
\subsubsection{Соберите готовый
журнал}\label{ux441ux43eux431ux435ux440ux438ux442ux435-ux433ux43eux442ux43eux432ux44bux439-ux436ux443ux440ux43dux430ux43b}}

Если вы уже перешли в \texttt{work/}, вернитесь в папку с материалами за
исходниками загрузчика и журнала, затем снова в рабочий каталог:

\begin{verbatim}
cd ..
cp -n boot16.asm work/boot16.asm
cp -n journal.asm work/journal.asm
cd work
\end{verbatim}

Общий загрузчик \texttt{boot16.asm} читает \textbf{17 секторов после
первого} и передаёт управление журналу по адресу \texttt{1000:0000}.
Соберите оба бинарных файла:

\begin{verbatim}
nasm -f bin boot16.asm -o boot.bin
nasm -f bin journal.asm -o journal.bin
wc -c boot.bin journal.bin
\end{verbatim}

\texttt{boot.bin} должен занимать ровно 512 байтов. Длину
\texttt{journal.bin} сравните с вместимостью 17 секторов по 512 байтов;
максимальное число байтов \textbf{посчитайте самостоятельно}. При
превышении нельзя продолжать сборку: текущий загрузчик не прочтёт
программу целиком.

Создайте образ пустой дискеты и поместите два бинарных файла на свои
места:

\begin{verbatim}
dd if=/dev/zero of=boot.img bs=512 count=2880
dd if=boot.bin of=boot.img bs=512 conv=notrunc
dd if=journal.bin of=boot.img bs=512 seek=1 conv=notrunc
wc -c boot.img
\end{verbatim}

\begin{longtable}[]{@{}
  >{\raggedright\arraybackslash}p{(\columnwidth - 2\tabcolsep) * \real{0.2700}}
  >{\raggedright\arraybackslash}p{(\columnwidth - 2\tabcolsep) * \real{0.6900}}@{}}
\toprule
\begin{minipage}[b]{\linewidth}\raggedright
Параметр \texttt{dd}
\end{minipage} & \begin{minipage}[b]{\linewidth}\raggedright
Что означает
\end{minipage} \\
\midrule
\endhead
\texttt{if=} / \texttt{of=} & Входной / выходной файл \\
\texttt{/dev/zero} & Источник нулевых байтов \\
\texttt{bs=512} & Размер блока в байтах \\
\texttt{count=2880} & Столько блоков по 512 байтов составляет дискету
1,44 MiB \\
\texttt{seek=1} & Отступить на один блок \emph{выходного} файла: начало
сектора 2 \\
\texttt{conv=notrunc} & Не усекать уже существующий образ при записи
фрагмента \\
\bottomrule
\end{longtable}

Теперь \textbf{один раз} создайте отдельный HDD для оценок. Не
повторяйте эту команду после того, как начали сохранять оценки: она
изменит длину образа, а для нового запуска нужен тот же самый файл с
данными.

\begin{verbatim}
truncate -s 1M grades.img
\end{verbatim}

Запуск полного журнала:

\begin{verbatim}
qemu-system-i386 \
  -machine pc -m 16M -boot order=a \
  -drive file=boot.img,format=raw,if=floppy \
  -drive file=grades.img,format=raw,if=ide,index=0,media=disk \
  -display gtk -serial stdio -monitor none -no-reboot
\end{verbatim}

Первый \texttt{-drive} --- дискета с программой. Второй \texttt{-drive}
--- \textbf{отдельный} RAW HDD. Параметры
\texttt{if=ide,index=0,media=disk} дают первый IDE-диск (для BIOS он
имеет номер \texttt{80h}). \texttt{-serial\ stdio} показывает
диагностику COM1 в вашем терминале; \texttt{-monitor\ none} выключает
командную консоль QEMU. Остальные флаги разобраны выше. Введите
\texttt{0}, стрелку вправо, \texttt{A}, нажмите \texttt{S}: ожидайте
\texttt{SAVED}. Завершите QEMU и запустите \textbf{той же командой}:
ожидайте \texttt{LOADED}, оценки на экране сохранились. Клавиша
\texttt{L} перечитывает диск. Если изменить исходник, пересоберите
\texttt{journal.bin} и повторите \textbf{все три команды \texttt{dd}}:
сначала получите чистый \texttt{boot.img}, затем поместите загрузчик и
программу. Это исключит остатки прежней, более длинной программы в
образе дискеты. \textbf{Не пересоздавайте \texttt{grades.img}:} полная
пересборка дискеты не меняет отдельный HDD.

\hypertarget{ux440ux430ux441ux441ux447ux438ux442ux430ux439ux442ux435-ux438-ux438ux437ux43cux435ux43dux438ux442ux435-ux43aux43eux43dux441ux442ux430ux43dux442ux44b}{%
\subsubsection{Рассчитайте и измените
константы}\label{ux440ux430ux441ux441ux447ux438ux442ux430ux439ux442ux435-ux438-ux438ux437ux43cux435ux43dux438ux442ux435-ux43aux43eux43dux441ux442ux430ux43dux442ux44b}}

Перед изменением запишите число, его основание и ожидаемое поведение.
Редактируйте только \texttt{journal.asm}. После каждой группы правок
снова выполните \texttt{nasm\ -f\ bin\ journal.asm\ -o\ journal.bin} и
три команды \texttt{dd} из раздела сборки журнала. Затем повторно
запустите QEMU с тем же диском оценок.

\hypertarget{ux432ux44bux432ux43eux434-ux446ux432ux435ux442-ux438-ux43aux43eux43eux440ux434ux438ux43dux430ux442ux44b}{%
\paragraph{Вывод: цвет и
координаты}\label{ux432ux44bux432ux43eux434-ux446ux432ux435ux442-ux438-ux43aux43eux43eux440ux434ux438ux43dux430ux442ux44b}}

В процедуре \texttt{vga\_text} найдите \texttt{mov\ ah,\ 0x0f}: это
атрибут каждого выводимого символа. Требуется \textbf{жёлтый текст на
синем фоне}, без мигания.

\begin{longtable}[]{@{}
  >{\raggedright\arraybackslash}p{(\columnwidth - 6\tabcolsep) * \real{0.1200}}
  >{\raggedright\arraybackslash}p{(\columnwidth - 6\tabcolsep) * \real{0.3600}}
  >{\raggedright\arraybackslash}p{(\columnwidth - 6\tabcolsep) * \real{0.1200}}
  >{\raggedright\arraybackslash}p{(\columnwidth - 6\tabcolsep) * \real{0.3600}}@{}}
\toprule
\begin{minipage}[b]{\linewidth}\raggedright
Код
\end{minipage} & \begin{minipage}[b]{\linewidth}\raggedright
Цвет
\end{minipage} & \begin{minipage}[b]{\linewidth}\raggedright
Код
\end{minipage} & \begin{minipage}[b]{\linewidth}\raggedright
Цвет
\end{minipage} \\
\midrule
\endhead
0 & чёрный & 8 & тёмно-серый \\
1 & синий & 9 & ярко-синий \\
2 & зелёный & A & ярко-зелёный \\
3 & бирюзовый & B & ярко-бирюзовый \\
4 & красный & C & ярко-красный \\
5 & пурпурный & D & ярко-пурпурный \\
6 & коричневый & E & жёлтый \\
7 & светло-серый & F & белый \\
\bottomrule
\end{longtable}

В атрибуте VGA бит 7 --- мигание (здесь 0), биты 6--4 --- фон (0--7),
биты 3--0 --- цвет символа (0--F). \textbf{Атрибут = 16×фон + цвет
символа}. Видеоячейка содержит байт символа и байт атрибута. Рассчитайте
новый операнд \texttt{mov\ ah,...}, выполните сборку и сравните цвет на
экране.

Для положения заголовка в \texttt{redraw} найдите
\texttt{mov\ bx,\ 0x0001} сразу перед \texttt{mov\ si,\ title\_text}.
Переместите заголовок на \textbf{строку 2, столбец 5} (нумерация с
нуля). \texttt{BH} задаёт строку, \texttt{BL} --- столбец;
\texttt{BX\ =\ 256×строка\ +\ столбец}. Смещение видеопамяти:
\texttt{2×(80×строка\ +\ столбец)} байтов от \texttt{B800:0000}.
Рассчитайте операнд \texttt{mov\ bx,...} и физический адрес первой
буквы. Заголовок не должен перекрывать подсказку в строке 1 и таблицу в
строке 3.

\hypertarget{ux432ux432ux43eux434-ux434ux440ux443ux433ux438ux435-ux43aux43bux430ux432ux438ux448ux438}{%
\paragraph{Ввод: другие
клавиши}\label{ux432ux432ux43eux434-ux434ux440ux443ux433ux438ux435-ux43aux43bux430ux432ux438ux448ux438}}

В \texttt{handle\_scan}, ветке \texttt{.normal}, найдите сравнения перед
\texttt{.save} и \texttt{.load}. Пусть \textbf{Q сохраняет, E загружает}
вместо S и L. Программа сравнивает скан-коды PS/2 set 1, а не ASCII. В
таблице указаны \textbf{десятичные} коды; переведите нужные
самостоятельно в hex:

\begin{longtable}[]{@{}lrrrrr@{}}
\toprule
Клавиша & S & L & Q & E & A \\
\midrule
\endhead
Make code, десятичный & 31 & 38 & 16 & 18 & 30 \\
\bottomrule
\end{longtable}

Исправьте также \texttt{help\_text} в конце файла. После перезапуска Q и
E должны работать, а старые S/L --- уже нет. Диагностику
\texttt{SAVED}/\texttt{LOADED} наблюдайте в терминале.

\hypertarget{ux432ux44bux447ux438ux441ux43bux435ux43dux438ux44f-ux43eux446ux435ux43dux43aux430-12}{%
\paragraph{Вычисления: оценка
12}\label{ux432ux44bux447ux438ux441ux43bux435ux43dux438ux44f-ux43eux446ux435ux43dux43aux430-12}}

Клавиша A теперь должна выставлять \textbf{12}, допустимый диапазон ---
0--12. Посмотрите \texttt{handle\_scan}, метку \texttt{.ten}, затем
\texttt{load\_grades}, метку \texttt{.check}: в первом месте находится
присваивание, во втором --- верхняя граница значения, принимаемого с
диска. Измените обе константы и \texttt{help\_text}.

Не меняйте делитель 10 в \texttt{vga\_two} и делитель 100 при выводе
среднего: они относятся к десятичной записи. Рассчитайте вручную среднее
для наборов \textbf{12, 0} и \textbf{12, 1, 1} с округлением до сотых.
Введите их и сравните экран с расчётом; сохраните 12, перезапустите и
убедитесь, что загрузка её принимает. Ноль --- выставленная оценка,
\texttt{FF} --- пропуск.

\hypertarget{ux434ux438ux441ux43a-ux434ux440ux443ux433ux43eux439-ux441ux435ux43aux442ux43eux440}{%
\paragraph{Диск: другой
сектор}\label{ux434ux438ux441ux43a-ux434ux440ux443ux433ux43eux439-ux441ux435ux43aux442ux43eux440}}

В пакете \texttt{dap} найдите \texttt{dq\ 1}, номер LBA. Требуется
сохранять с \textbf{байта 4096} образа HDD. Сектор имеет 512 байтов:
самостоятельно вычислите \texttt{LBA\ =\ 4096\ /\ 512}, затем поменяйте
значение \texttt{dq}.

Перед опытом сохраните прежний диск как отдельный файл и собирайте
программу описанными выше командами:

\begin{verbatim}
cp -n grades.img grades-lab.img
qemu-system-i386 \
  -machine pc -m 16M -boot order=a \
  -drive file=boot.img,format=raw,if=floppy \
  -drive \
    file=grades-lab.img,format=raw,if=ide,index=0,media=disk \
  -display gtk -serial stdio -monitor none -no-reboot
\end{verbatim}

Первый запуск с новым LBA может сообщить \texttt{EMPTY}: старое
сохранение остаётся в другом секторе. Выставьте 12 и 0, сохраните Q и
ещё раз запустите QEMU с \textbf{той же копией}. Должно появиться
\texttt{LOADED}. Подставьте своё рассчитанное смещение и прочитайте
сигнатуру:

\begin{verbatim}
OFFSET=ваше_десятичное_смещение
od -An -tc -j "$OFFSET" -N 4 grades-lab.img
\end{verbatim}

\texttt{-tc} показывает байты как символы; \texttt{-j} пропускает нужное
число байтов, \texttt{-N\ 4} читает четыре байта. Убедитесь также, что
старый сектор в исходном \texttt{grades.img} не изменился.

\hypertarget{ux43eux431ux44aux44fux441ux43dux438ux442ux435-ux43fux443ux442ux44c-ux43eux446ux435ux43dux43aux438}{%
\subsubsection{Объясните путь
оценки}\label{ux43eux431ux44aux44fux441ux43dux438ux442ux435-ux43fux443ux442ux44c-ux43eux446ux435ux43dux43aux438}}

Проследите путь оценки: клавиша → скан-код → числовой байт в RAM → два
символа и атрибуты VGA → буфер \texttt{sector} в RAM → сектор HDD. Чем
отличаются \texttt{1000:0000} (код), \texttt{B800:...} (видеопамять) и
\texttt{1000:sector} (буфер для диска)?

\hypertarget{ux441ux430ux43cux43eux441ux442ux43eux44fux442ux435ux43bux44cux43dux430ux44f-ux440ux430ux431ux43eux442ux430-ux434ux43eux43cux430}{%
\subsubsection{Самостоятельная работа
дома}\label{ux441ux430ux43cux43eux441ux442ux43eux44fux442ux435ux43bux44cux43dux430ux44f-ux440ux430ux431ux43eux442ux430-ux434ux43eux43cux430}}

\begin{enumerate}
\def\labelenumi{\arabic{enumi}.}
\item
  \textbf{C и GDB: файл и процесс.} Из \texttt{work/} выполните команды
  в указанном порядке:

\begin{verbatim}
 gcc -std=c11 -Wall -Wextra -Werror -g -O0 \
   ../memory.c -o memory
 file memory
 ./memory
 echo $?
 gdb -q ./memory
\end{verbatim}

  \texttt{-std=c11} выбирает стандарт языка;
  \texttt{-Wall\ -Wextra\ -Werror} включает диагностику и превращает
  предупреждения в ошибки; \texttt{-g} добавляет отладочные сведения;
  \texttt{-O0} отключает оптимизации для пошагового разбора;
  \texttt{-o\ memory} задаёт имя программы. \texttt{file\ memory} должен
  узнать исполняемый файл ELF для Linux x86-64; детали описания зависят
  от версии GCC. \texttt{./memory} просит Linux запустить файл:
  появляется процесс со своим виртуальным адресным пространством. Сразу
  после завершения программы \texttt{echo\ \$?} показывает её код
  завершения --- здесь ожидается \texttt{0}, если проверки в
  \texttt{memory.c} прошли. \texttt{gdb\ -q} убирает вводный баннер
  отладчика. \textbf{Вводите команды GDB вручную}:

\begin{verbatim}
break show_memory
run
print sizeof(char)
print sizeof(int)
print sizeof(long)
print *x
x/4xb x
x/3dw numbers
continue
quit
\end{verbatim}

  \texttt{break} ставит остановку; \texttt{run} запускает программу;
  \texttt{print} вычисляет выражение; \texttt{x/4xb} читает четыре
  единичных байта в hex, \texttt{x/3dw} --- три 32-битных слова в
  десятичном виде; \texttt{continue} продолжает работу. Объясните шаги
  адресов массивов \texttt{char/int/long} и little-endian представление
  \texttt{0x12345678}. Адрес \texttt{x} в GDB --- виртуальный адрес
  запущенного процесса Linux, а не физический адрес журнала в QEMU; при
  других запусках виртуальный адрес может измениться.

  \textbf{Вопросы к лекции №5.} Что здесь является исходным файлом, что
  --- исполняемым файлом на диске, а когда существует процесс? Почему
  \texttt{./memory} запускается в Linux, а \texttt{hello.bin}
  исполняется только после загрузки BIOS в QEMU? Кто предоставляет
  программе адресное пространство и возвращает оболочке код завершения?
  На следующей теме разберём загрузку программы, процессы и услуги
  операционной системы.
\item
  \textbf{Цвета и адреса.} Вычислите по таблице ещё два атрибута:
  «ярко-зелёный на чёрном» и «белый на красном». Для двух свободных
  позиций экрана укажите \texttt{BX}, смещение VGA и физический адрес.
\item
  \textbf{Ввод и арифметика.} Запишите перевод десятичных скан-кодов Q/E
  в hex, вычислите средние для 12,0 и 12,1,1; объясните разницу между
  числом 12 в \texttt{grades} и текстом на экране. Почему
  \texttt{vga\_two} по-прежнему делит на 10?
\item
  \textbf{Сохранение.} На копии \texttt{grades-lab.img} исследуйте
  \texttt{GR16}, четыре байта заголовка, оценки и контрольную сумму. Для
  опыта с повреждением создайте \textbf{дополнительную} копию и измените
  один байт оценки, оставив сумму прежней:

\begin{verbatim}
cp -n grades-lab.img grades-corrupt.img
printf '\001' | dd of=grades-corrupt.img bs=1 \
  seek="$((OFFSET + 8))" conv=notrunc
\end{verbatim}

  Здесь
  \texttt{printf\ \textquotesingle{}\textbackslash{}001\textquotesingle{}}
  передаёт один байт со значением 1, \texttt{bs=1} задаёт байтовое
  смещение, \texttt{seek} начинает запись по вычисленному адресу первого
  поля, \texttt{conv=notrunc} сохраняет остальной образ. Сначала
  убедитесь, что на этом месте была оценка 12: иначе выберите другой
  байт. Запустите журнал с \texttt{grades-corrupt.img} вместо
  \texttt{grades-lab.img}: наблюдайте \texttt{BAD\_CHECKSUM}. После
  ввода новой оценки нажмите E и убедитесь по экрану, что плохая
  загрузка её не затёрла. При повторном опыте выберите \emph{новое} имя
  копии.
\item
  \textbf{Протокол и отчёт.} Представьте свои вычисления и изменённый
  \texttt{journal.asm}. Приложите результат
  \texttt{diff\ -u\ ../journal.asm\ journal.asm}: исходный файл лежит в
  папке материалов, ваша версия --- в рабочем каталоге. Для каждого
  опыта составьте таблицу «ожидалось / наблюдалось»: первый запуск,
  цвет, позиция, клавиши, оба средних, сохранение, загрузка и
  повреждённый диск. Описывайте проверку своими наблюдениями. Для
  структуры можно посмотреть «Пример отчёта: сумма двух чисел» в
  дополнительных материалах; его числовые результаты к журналу не
  относятся. Отправьте преподавателю свой отчёт в формате PDF или HTML,
  сохранив исходники и образы отдельно для демонстрации.
\end{enumerate}

\hypertarget{ux444ux43eux440ux43cux430-ux43eux442ux447ux451ux442ux430-ux434ux43bux44f-ux43aux430ux436ux434ux43eux433ux43e-ux438ux437ux43cux435ux43dux435ux43dux438ux44f}{%
\paragraph{Форма отчёта для каждого
изменения}\label{ux444ux43eux440ux43cux430-ux43eux442ux447ux451ux442ux430-ux434ux43bux44f-ux43aux430ux436ux434ux43eux433ux43e-ux438ux437ux43cux435ux43dux435ux43dux438ux44f}}

\begin{verbatim}
Задание и место в исходнике (метка, инструкция):
Мой расчёт (с основанием системы счисления и единицами):
Старая и новая константы:
Команды, которыми я собрал и запустил программу:
Ожидаемый результат:
Наблюдаемый результат (экран / терминал / байты образа):
Объяснение совпадения или расхождения:
\end{verbatim}
